\chapter{Methodology and Experiments}  
\label{chap:methodologyExperiment}

This chapter describes the experimental design, setup, and execution for capturing, analyzing, and understanding the traffic generated by the Ray-Ban Meta glasses and its companion app. It also addresses the limitations encountered during the process.

\section{Experiment Design}
The experiments aim to identify the privacy risks associated with the Ray-Ban Meta glasses and its companion app. The goal of the study includes identifying the types of data generated by the glasses and app, examining the encryption practices, determining the destinations of the transmitted data (geological and corporational), and interactions with third parties, such as advertisers and trackers, and assessing whether these transmissions include personal information about users or bystanders. An additional focus is paid to the permissions requested by the companion app and third-party libraries used by the app for its potential privacy risks as it may expand the data sharing beyond user awareness. 

The experiments were designed to prioritize interactions between the Ray-Ban Meta glasses and the companion app on the phone, with efforts to minimize external interference and ensure the capture of relevant traffic. However, it is important to acknowledge that this controlled environment may not fully replicate the complexity of real-world scenarios. 

The glasses rely on both network and Bluetooth connectivity, and upon that, each communication serves primarily different purposes. Therefore, the final experiment design has separate captures and analysis methods for the different communications. Finally, all captures were also made with bystander privacy in mind, where the photos and video recordings were taken facing a white wall with no third-party voice other than the audio streaming from Spotify captured during any voice recordings. 

The research initially aimed to analyze the decrypted traffic for sensitive data such as raw readings. A man-in-the-middle (MitM) setup using mitmproxy was used to intercept and decrypt HTTPS streams for network traffic. As outlined in \cite{fereidouni2023iotmaninthemiddleattacks}, IoT devices often lack solid security measures to protect transmitted data. However, the Meta View app uses certificate pinning, which limits this approach. Other tools, like Frida and Objection, were then used to try and bypass the pinning, but the strong certificate pinning remained. These technical limitations prevented full decryption, resulting in the analysis focusing on metadata, traffic patterns, and destinations to assess privacy risks.


\section{Experimental Setup}
The environmental setup used to conduct the experiments are detailed in this section. It outlines the hardware and software components involved, including the testbed and tools used to capture the traffic generated by the glasses and companion app.

\subsection*{Tools Used}
Wireshark \cite{Wireshark}, an app that captures data packets in real time, was the primary tool for monitoring network traffic and Bluetooth communications. It provides a packet-level view of the data exchanged to examine the traffic patterns, protocols, and information transmitted between the glasses and the phone. An ALFA network adapter (Long-Range USB Adapter of model AWUS036NHA) was used to create a Wi-Fi hotspot, allowing the laptop to act as an intermediary to route network traffic and capture it for analysis. Classic Bluetooth captures were obtained using the Ubertooth One \cite{UbertoothOne}, a sniffer device with firmware version 2024.10. Meanwhile, the nRF52840 Dongle \cite{nRF52840}, another sniffer device, was used to monitor BLE traffic. A man-in-the-middle proxy (mitmproxy \cite{mitmproxy}) was used to intercept and decrypt network traffic to analyze requests, responses, and payloads, though certificate pinning restricted much analysis. An additional phone was used as a recording device to have an auditory record of the experiments. 

For the analysis of the captured traffic, visualizations were generated post-capture using custom Python scripts. The analysis of the companion app was done using Mobile Security Framework (MobSF) \cite{MobSF}, an automated security assessment tool for mobile applications. It identifies permissions and third-party integrations embedded within the Android Application Package(APK), a file format that is used to install an application on an Android operating system.



\subsection*{Testbed}
The primary components in the testbed for the experiments included a laptop running Ubuntu Linux version 24.04, a hotspot created using an ALFA network adapter connected to the laptop, the Ray-Ban Meta glasses, and a Samsung S24 smartphone running on Android 14 as the main device. The hotspot only had the Samsung S24 connected to isolate the traffic captured through Wireshark. For Bluetooth communications, a sniffer was plugged into the laptop via USB to capture the traffic. The Bluetooth sniffer nRF52840 Dongle was used for BLE experiments, and the Ubertooth One was used for Classic Bluetooth experiments. For network traffic, the capture option for the experiments was set to the hotspot, and for BLE traffic, the interface was set to the dongle to capture traffic on those interfaces.

Three phones were used throughout the thesis. During the earlier phases, two unsupported phones, a Xiaomi Redmi Note 8 and a Xiaomi POCO X5 5G, were used to explore device functionalities, test preliminary setups, and attempt traffic decryption. However, due to the instability and compatibility problems, a third phone was introduced, a Samsung S24, which was an officially supported phone model. All final experiments were performed on this phone, ensuring consistency and comparability in the results.




\section{Data Collection}
This section covers the categorization of traffic data generated by the Ray-Ban Meta glasses and companion app. The collected data was split into two main categories: network traffic and Bluetooth traffic. The data in each main category was then further divided based on the need for analysis and experimentation, as seen in Figure \ref{fig:traffic_categories}.

\begin{figure}
    \centering
    \includegraphics[width=0.8\textwidth]{Figures/ME_traffic_categories.png}
    \caption{Categorization of the types of Traffic generated by the Ray-Ban Meta glasses and the companion app}
    \label{fig:traffic_categories}
\end{figure}

\subsection*{Network Traffic}
For Network Traffic, data was separated into Idle and Active traffic. Idle traffic involves no direct user interaction and was captured while the glasses remained idle for an extended period of time. It is a controlled test that helps identify persistent background data flows. The second type is active traffic, which refers to the traffic captured whilst interacting with the glasses and companion app.

From initial experimentation, it was noted that the Meta View app does not allow pairing when there is no internet connection. Therefore, active traffic is divided into two subcategories: traffic from daily interactions, such as taking a photo or asking the Meta AI about the time, and initial pairing traffic, which is exchanged between the phone and glasses when establishing the connection between the phone and glasses. Furthermore, the preliminary overview and hands-on usage of the glasses identified two types of commonly performed daily interactions. Thus, daily interactions are further grouped into two types of commands: touch-based and voice-based commands. Touch commands execute common user interactions through manual inputs, while voice commands are carried out through the AI voice assistant, allowing for a broader range of interactions during experimentation. 

\subsection*{Bluetooth Traffic}
Bluetooth Traffic was split based on its type. The first is Bluetooth Low Energy (BLE), which is designed for low-bandwidth, energy-efficient communication and is typically used for the initial pairing of devices. The second is Classic Bluetooth, which supports higher-bandwidth communication and is used for continuous data streaming, such as audio playback. 

Bluetooth traffic is event-driven, meaning that the traffic captures are primarily during active interactions with the glasses and phone. Unlike network traffic, which may include background traffic even when idle, Bluetooth traffic is less likely to have meaningful data when kept idle; therefore, there is no idle traffic experiment for Bluetooth. Nonetheless, the data for both types of Bluetooth traffic was divided into user interaction and initial pairing of the devices, of which user interactions are then grouped based on the type of interaction, touch-based or voice-based commands.


\section{Experimental Procedure}
This section describes the process of setting up the environment and capturing the generated traffic. The experiments were repeated multiple times to ensure better reliability. Each round of experiment capture is referred to as a trial in the analyses. The following subsections first introduce the experimental procedure for different traffic types (network and Bluetooth) before discussing further details for each experiment in their respective subsections. 

\subsection*{Network Traffic}
The environment was set up with only the phone connected to the hotspot. The Ray-Ban Meta glasses were then powered off, and the companion app was forced to quit to kill its process. The experiment began by simultaneously starting a Wireshark capture on the laptop and the voice recording on the extra phone. The Wireshark was set only to capture traffic going through the hotspot. The procedure for the remaining process of the different experiments, including active interactions and initial pairing, will be discussed in the following subsections. The experiments for active interactions with the network traffic were repeated five times each.

\subsection*{Bluetooth Traffic}
The only difference between the setup for Classic Bluetooth and the BLE was the sniffer used to capture the nearby Bluetooth traffic. Classic Bluetooth used the Ubertooth one, whilst the nRF52840 was used for BLE. The respective Bluetooth sniffers were plugged into the laptop via USB. The Ray-Ban Meta glasses were powered off and the companion app was force quitted. The Ubertooth one is started via a command on the terminal, whereas the nRF52840 Dongle was started by selecting its interface in Wireshark for the capture. The experiment began by simultaneously starting the Bluetooth capture using the sniffers and voice recording. The following subsections will discuss the procedure for the remaining process of the different experiments for active interactions and initial pairing. All experiments for Bluetooth traffic were repeated twice.

\subsection*{Touch-based commands}
The touch-based commands mainly focused on Media Capture and Music Playback. After the setup of the environment, a predefined list of commands was executed manually. This list can be found in \cref{tab:touch_commands}. 
The execution order started with Media capture, which was then imported onto the phone. Next, the commands for Music Playback were performed, starting with playing music. Up until this interaction, each command was given at thirty-second intervals. The time interval allows for sufficient time for the traffic to stabilize between each interaction. Afterwards, the rest of the Music Playback commands were given in fifteen-second intervals. The duration of each touch-command experiment was roughly five minutes.

\subsection*{Voice-based commands}
Voice commands have a wider range of functionalities than touch commands. The different functionalities of voice commands include General Questions, Music Playback, Timers, Media Capture, Messaging, and Phone Calls. Similarly, before the start of the experiment, the environment was set up. The specific set of commands executed can be found in \cref{tab:voice_commands}. All voice commands were performed in thirty-second intervals, ensuring consistency. Thereby making it easier to compare and observe patterns. An automated voice recording delivered the initial series of voice commands.

The specific sequence of functionalities that used the automated recording includes General Questions, Music Playback, Timers, and Media Capture. The remaining commands of the experiment were performed manually with the same thirty-second interval per command. The transfer of captured media from the glasses to the phone was the first manual command executed, followed by the commands for Messaging and Phone Calls. The duration of each voice command experiment was roughly eleven and a half minutes.

\subsection*{Initial Pairing}
Pairing Sequence handles traffic during the initial pairing of the glasses and the companion app. This experiment had one more prerequisite in the setup before starting the experiment. The glasses needed to be unpaired and forgotten on the phone in Bluetooth settings. The experiment then began when the glasses were powered on and the app opened. The interaction then continues with the glasses being paired with the phone. The glasses are then unpaired and forgotten again before re-pairing. Each pairing round took roughly one minute. For this experiment, the unpairing and re-pairing actions were repeated five times but in the same traffic capture. This means that the five repetitions of pairing between the devices were done in one go, whereas the experiments for the touch-based and voice-based commands had five separate captures, one for each repetition. The duration of the initial pairing experiment was around five and a half minutes.

\subsection*{Idle Traffic}
A single capture of approximately eight hours took place throughout the night to observe whether the glasses communicated with the Meta servers or other endpoints. No interactions were done on the phone or the glasses. The experiment started with the phone and glasses being powered on and left in proximity while running the Wireshark capture.


  
\section{Analysis Methods}
The captured traffic was analyzed offline after all experiments were completed. This section describes the process and methods used to analyze each type of traffic. Wireshark was used primarily for packet inspection, while the Python scripts generated visualizations for the traffic volume, packet sizes, protocol usage, and destination IP mapping for each experiment. The captured data was exported as a CSV file, aggregated in Python, saved as a CSV file, and subsequently visualized using standard Python libraries for analysis.


\subsection*{App Analysis}
The companion app for the Ray-Ban Meta glasses was analyzed using the static Mobile Security Framework (MobSF), a tool specifically designed for analyzing mobile applications. It decompiles Android applications' APK files to inspect their source code, resources, and configuration files. This tool detects hardcoded credentials, requested permissions, and integration of third-party libraries. For the app analysis, the components found were divided into four types. The first are first parties, which refers to the components or elements provided by Meta. Support parties refer to those that come from the same parent organization, Meta, and support functionalities and features of the app but are external to the app itself (e.g., WhatsApp). Finally, third parties are those that are not provided by Meta Platforms but instead by an external, independent organization (e.g., Google). The fourth is platform-provided, which are components provided as part of the Android development framework or operating system (e.g., AndroidX).

\subsection*{Network Traffic}
\label{subsec:networkTraffic}
In general, the network traffic captured is divided into four different parties in the analysis. The first is private traffic, which includes local broadcasts or traffic directed to the Samsung S24 phone. The second is first-party traffic, which refers to the destination addresses owned by Meta and is directly related to the app (e.g., the Meta servers). Support-party traffic refers to those that are sent to addresses that are owned by Meta but are external to the app (e.g., WhatsApp). Finally, third-party traffic refers to those that are not owned by Meta Platforms and are instead routed to an external, independent organization (e.g., Google). In order to distinguish whether traffic went to private, first-party, support-party, or third-party domains, the packets were grouped based on IP ownership and DNS resolution and then categorized accordingly.

Network traffic was also grouped based on its readability to understand the extent of information exposure better. The protocol was found to be an effective distinguisher for assessing readability, making it the primary criterion for grouping the data. A packet is considered readable if most of the packet content is human-readable. It is considered partially readable if a minor part of the payload can be understood. Finally, it is considered unreadable if only the metadata is exposed and the payload is completely not human-readable.

The captured traffic showed some third-party hostnames that were labeled as Google. These are part of Google's global network, including googleapis.com. At the same time, some other hostnames were labeled as Google Cloud, which corresponds to the use of Google Cloud Platforms (GCP). These are often used by other applications and services for hosting or storage. This brings attention to a crucial differentiation. While some IP addresses may point to Google servers, the actual organization being communicated with may not be Google itself but rather a third-party service using Google's infrastructure. These addresses were reverse DNS resolved and analyzed within the captures to determine their associated services. In the captures, the main application found relying on Google Cloud for hosting and delivering services was Spotify. 


\subsection*{Bluetooth Low Energy Traffic}
The BLE captures revealed only a few MAC addresses, which is a positive indicator for the isolated environment created for the experiment. However, none matched the MAC address of the glasses or the phone. Assuming that the BLE communication between the phone and glasses was captured, this suggests that the devices are using randomized MAC addresses for communication. In order to find which addresses belong to our devices, the Advertising Data for each packet was extracted. In Advertising Data, there is a field in the packet called Manufacturer Specific, which includes a company ID that provides the manufacturer of the device sending that packet. The Company ID was used as a filter, and the following two IDs, Samsung Electronics Co. Ltd (0x0075) and Meta Platforms, Inc. (0x01ab), were found. These are the two manufacturing companies of the devices involved in the experiments. These IDs were used as a filter, for example, btcommon.eir\_ad.entry.company\_id == 0x0075 to determine the number of MAC addresses involved for each ID. Only one MAC address was associated with each company, 5f:64:8e:90:8d:44 for the glasses and SamsungElect\_98:02:66 for the phone.


\subsection*{Classic Bluetooth Traffic}
An overview of captures shows that packets lack a source and destination, which are typically used to identify communication endpoints. Therefore, in order to identify which Classic Bluetooth communication was from the glasses or the phone, the Lower Address Part (LAP) was extracted from the Bluetooth Pseudoheader in each packet. Bluetooth device addresses are unique identifiers assigned to a device by its manufacturer, typically 48 bits in length. The addresses consist of a 16-bit Non-significant Address Part (NAP) and an 8-bit Upper Address Part (UAP); these bits usually make up the identifier assigned to the organization. The last 24 bits are the Lower Address Part (LAP), a unique identifier the manufacturer assigns to a device. 

The glasses are powered on at minute 0.5 for each experiment conducted, meaning that any traffic occurring before that time could not have originated from the glasses. For instance, Figure \ref{fig:avg_packet_count} displays the packet count for all addresses captured during the touch-based commands experiment. The address 0x0021c4f4, represented with the blue line, was captured within the first thirty seconds of the experiment. Since the glasses were not yet powered on, this LAP could not belong to them and was therefore eliminated as a candidate.

\begin{figure}[H]
    \centering
    \includegraphics[width=0.9\textwidth]{Figures/CT_avg_packet_count_all_traffic.pdf}
    \caption{Average packet count for all captured addresses in the touch-based command experiment.}
    \label{fig:avg_packet_count}
\end{figure}

Another method used to filter the addresses involved the signal strength. Given that the glasses and phone were placed right next to the sniffer, they should have high signal strength. However, there was no RSSI value (Received Signal Strength Indicator) in the capture, as it is not provided by the Ubertooth one. This was a hardware limitation, which otherwise would have indicated the strength of the received signal at the receiver's end. Instead, a signal power (SP) value was available, measuring a relative signal level on a channel, which could be used to infer signal quality in place of RSSI. The signal power is examined to determine what addresses would make the most sense. The average signal strength for the experiment environment should be below -60 dBM due to the placement of the sniffer and the devices. The average signal strength for the LAPs, after filtering out those found within the first thirty seconds of the experiment, was extracted and depicted in Figure \ref{fig:avg_sp}, which showed only two possible candidates for the devices. 

In order to check which of the LAPs belonged to the glasses, the traffic for each LAP over time was tracked to determine whether it aligned with the interaction timings where a command was issued. Two LAPs were found to align well with the commands and each other, those being 0x00347a8c with an average SP of -23 and 0x002ff7a8 with an average SP of -38 dBm, as seen in Figure \ref{fig:avg_sp}. These LAPs were then cross-referenced between the different experiments and the trials within each experiment using, for example, btbredr\_rf.lower\_address\_part == 0x00347a8c as a filter in Wireshark. There was a possibility that the devices used Bluetooth address randomization, which would allow a device to regularly change its Bluetooth addresses to prevent long-term tracking by others. Fortunately, all classic Bluetooth experiments performed were done in close periods to one another, and both LAPs were present in all captures.

\begin{figure}[H]
    \centering
    \includegraphics[width=0.9\textwidth]{Figures/CT_avg_sp_all_after_30_traffic.pdf}
    \caption{Average signal power (SP) over time for remaining addresses in the touch-based commands experiment after filtering.}
    \label{fig:avg_sp}
\end{figure}
    
The next question was to identify which LAP belongs to the phone and which to the glasses. However, note that the correctness of this assumption does not affect the analysis. Here, the traffic frequency for each LAP was used, which identified one LAP to be active only for five seconds, roughly when the glasses are paired, and the other LAP to be active a bit before the first LAP and continue to be active throughout the experiment.

Classic Bluetooth typically has a master-worker relationship in which one device coordinates the communication, and the other device responds to its commands and requests; therefore, the second identified LAP would be the master device. A phone would usually be considered a master device due to its controls and ability to connect to multiple devices. A worker device can only connect to its master device and cannot communicate with any other worker devices. However, this is only a general premise; therefore, the sequence of packets during pairing was examined to identify which LAP was the first to begin communication. 

Since the master scans the surrounding devices and selects the worker device to connect to, it would make sense that the master device is the one to initiate communication. The captures confirmed that the second LAP is the master device, which is believed to be the phone since the phone initiates the pairing with the glasses. For certain operations that might need the glasses to be the master, it is important to know the design of Classic Bluetooth, where the master controls the communication timing and surveys the worker for data. Even if the glasses initiate commands or stream data, the phone could technically still be the master that is coordinating all transmissions. 

Classic Bluetooth also allows for dynamic role swapping, which allows the worker to have temporary control and may explain some commands where the glasses appear to be the ones giving the command. Therefore, it is the assumption that the phone is of the address 0x002ff7a8 and is the master device controlling the link and the glasses, originally on the address 0x00347a8c, to use the link of the phone for further communication. 



\section{Limitations}
Several limitations affected the scope and progression of this thesis. First, the unsupported phone models proved to be much more of a problem than initially predicted. Second, the process of capturing the Bluetooth packets proved to be challenging. Lastly, the inability to decrypt the encrypted traffic restricted the analysis to mainly metadata and destination mappings. The following sections go into detail about the limitations of each aspect.

\subsection*{Phone}
The initial experiments were performed with the companion app on a Redmi Note 8 running Android 14. This phone was able to perform most functions. However, live streaming did not work as the app would crash when trying to live stream with glasses. The Bluetooth connection between the glasses and the phone was also relatively unstable, with the phone having problems finding the glasses. This phone model also had constant problems with transferring media from the glasses to the phone. The transfer process involved the glasses establishing their own Wi-FI network, acting as a Wi-Fi access point to which the phone connects for the duration of the transfer. As a result, fixes, including unpairing and re-paring the glasses, factory resetting the glasses, and deleting and reinstalling the app, were common occurrences.

The decryption efforts on network traffic were made on a POCO X5 5G phone running Android 14 with the security update dated October 2024. The limitation of this phone was much more noticeable as a simultaneous connection of Bluetooth and the app was not possible. As a result, the user can choose to either pair the glasses with Bluetooth in Settings, which allows music playback, or connect to the app, which allows for media transfer, but not both at the same time. This limitation was far more detrimental to the ability to perform the experiments. The Bluetooth connection between the glasses and the phone was also relatively unstable as the glasses attempted to connect via Bluetooth but failed repeatedly due to the limitation that prevents simultaneous pairing of the devices over both Bluetooth and the network.

The third phone used was a Samsung S24 on Android 14. This phone is one of the phone models officially supported by the Meta View app and had none of the previous connection problems, which allowed for a much smoother experience. This phone model was selected to be used for the final experiments.

The final limitation of the phone was the geographical location. The Ray-Ban Meta glasses are available and supported in multiple countries, including the US, UK, Germany, and others. However, Switzerland is not one of those countries. Therefore, some features of the Meta AI were restricted. These features include asking the Meta AI to translate text, set reminders, ask the AI to look at an object where it would try to identify it or ask the AI to scan QR codes or call a number placed within the camera view of the glasses. This limited the number of interactions that could be included in the experiment design.

\subsection*{Bluetooth Sniffer}
Multiple sniffers were considered in the initial experimentation for Bluetooth sniffing. However, some could not reliably capture the communications, resulting in the decision to use the Ubertooth One and the nRF52840 Dongle for the final experiments. These devices were able to sniff the packets exchanged between the glasses and the phone reliably. 

Initially, the Micro:bit V2.20/21 \cite{Microbit}, a programmable microcontroller board, flashed with a sniffer firmware, btlejack \cite{btlejack}, was used to capture BLE packets. However, the Micro:bit was unsuitable for this experiment as it could not reliably sniff the connection between the phone and the glasses. The way for the Micro:bit to target a connection requires the Micro:bit to first sniff for existing connections and retrieve the access address of active BLE devices, which is then output into the terminal as a list. Next, it is necessary to determine which access address corresponds to the BLE device and then use that access address to target the connection between the devices. The problem was that the sniffed addresses had RSSI values that were too low to be from to the glasses. Given that the glasses were positioned directly beside the sniffer, such low signal strength was unlikely. Furthermore, interactions with the glasses did not affect the packet count for these addresses. The Micro:bit, in general, was not able to sniff many connections during initial experimentations.

Another sniffer tried was the Adafruit Bluefruit nRF51822 \cite{BluefruitSniffer}, a board that can monitor BLE traffic between devices when flashed with sniffer firmware \cite{AdafruitBLESnifferPython}. However, the Adafruit was also found to be unsuitable for this experiment as it was also unable to sniff the connection between the phone and the glasses. The way for the Adafruit to target a connection requires it to first sniff for existing connections and retrieve the MAC address of active BLE devices, which is then output into the terminal as a list alongside an incrementing number assigned to each MAC address. Next, the number associated with the device's MAC address must be used to target the connection between the devices. When testing the Adafruit, a list of at least 20 different MAC addresses was presented, which was more than the Micro:bits. However, the MAC address of the glasses does not appear in the list.

The third sniffer considered for sniffing BLE traffic was the nRF52840 Dongle. It is a USB device that can monitor BLE traffic and integrates with Wireshark directly, allowing for better real-time, on-screen visualization of the captured packets. Although it required more setup for the flashing, requiring a programmer app on the desktop to flash the firmware \cite{nRF52840firmware}. The nRF52840 Dongle targets connections by scanning for nearby BLE devices, identifying their MAC addresses, and allowing the user to select a target connection to monitor. However, this feature was not needed as the on-screen visualization on Wireshark was enough to identify the Ray-Ban Meta glasses and the phone. During experimentation with the device, the nRF52840 performed much better than the other devices in sniffing BLE traffic, as it was able to consistently detect the glasses and phone during pairing. These advantages made it the final choice for the BLE sniffing experiment.

The fourth and final sniffer considered was the Ubertooth One, a USB device that is able to scan, capture, and monitor Bluetooth traffic across channels. Unlike the previous sniffers, it comes pre-flashed with firmware, allowing it to function as a Bluetooth and BLE sniffer right out of the box. This means that there is no need to flash additional firmware. It can also sniff Classic Bluetooth, which none of the previous devices supports. This means that it was the only choice for Classic Bluetooth sniffing. Initially, the Ubertooth One was chosen for BLE sniffing as well. The reason was that its process did not require sniffing existing connections in order to target a specific connection. For the Ubertooth One, only the MAC address of the device was needed to target a device. However, the resulting capture would sometimes have random time jumps during traffic capture, therefore resulting in the nRF52840 Dongle being used for BLE sniffing.